\chapter{构建、测试和调试}

\section{问题入口：一份源文件走不远}

单文件程序适合学习语法，但真实项目很快需要多个源文件、测试和构建脚本。假设项目有一个统计库：

\begin{lstlisting}[caption={最小项目结构}]
score_tool/
  CMakeLists.txt
  include/score/summary.hpp
  src/summary.cpp
  tests/summary_tests.cpp
\end{lstlisting}

目录结构表达所有权：公开头文件在 \code{include}，实现放 \code{src}，测试放 \code{tests}。不要把所有文件堆在一个目录里。

\section{CMake 的基本边界}

\begin{lstlisting}[language=CMake,caption={库和测试目标分开}]
cmake_minimum_required(VERSION 3.20)
project(score_tool LANGUAGES CXX)

add_library(score_summary src/summary.cpp)
target_compile_features(score_summary PUBLIC cxx_std_20)
target_include_directories(score_summary PUBLIC include)

add_executable(summary_tests tests/summary_tests.cpp)
target_link_libraries(summary_tests PRIVATE score_summary)
enable_testing()
add_test(NAME summary_tests COMMAND summary_tests)
\end{lstlisting}

目标是 CMake 的核心。不要把编译选项、include 路径和库依赖全局乱撒。谁需要依赖，谁声明依赖。

\section{测试不是最后补}

统计函数天然适合测试：

\begin{lstlisting}[language=C++,caption={不用测试框架也能先写断言}]
#include <cassert>
#include <vector>

int main()
{
    const std::vector<int> scores{90, 80, 100};
    const auto summary = summarize_scores(scores);
    assert(summary.has_value());
    assert(summary->count == 3);
    assert(summary->sum == 270);
    assert(summary->min_score == 80);
    assert(summary->max_score == 100);
}
\end{lstlisting}

真实项目可以使用 Catch2、GoogleTest 或 doctest。框架不是重点，重点是测试要覆盖边界：空输入、单元素、多元素、极值和错误路径。

\section{调试先缩小问题}

调试不等于到处打印。更可靠的流程是：

\begin{enumerate}
  \item 复现问题，固定输入。
  \item 写一个最小失败测试。
  \item 确认失败位置：输入解析、核心逻辑、输出格式还是环境。
  \item 修正后保留测试，防止回归。
\end{enumerate}

调试器可以观察变量、调用栈和线程状态；日志可以保留运行历史；断言可以检查不变量。三者用途不同。

\begin{lstlisting}[language=C++,caption={断言用于检查内部不变量}]
#include <cassert>

void push_value(std::vector<int>& values, int value)
{
    const std::size_t old_size = values.size();
    values.push_back(value);
    assert(values.size() == old_size + 1);
}
\end{lstlisting}

断言不是用户错误处理。它适合检查“代码内部本应永远成立”的条件。

\section{本章边界检查}

一个能长期维护的 \cpp 项目至少要有：

\begin{itemize}
  \item 清晰目标：库、程序、测试分别是什么。
  \item 目标级 include 路径和依赖。
  \item 可重复运行的测试命令。
  \item Debug 和 Release 两种构建。
  \item 编译警告作为质量输入。
\end{itemize}

\begin{keyidea}
工程化不是等项目变大后才需要。构建边界、测试入口和调试习惯越早建立，后面修改成本越低。
\end{keyidea}

\section{从一条编译命令演进到 CMake}

单文件时可以直接编译：

\begin{lstlisting}[language=bash,caption={单文件编译}]
g++ -std=c++20 -Wall -Wextra -pedantic main.cpp -o score_tool
\end{lstlisting}

当项目拆成头文件、源文件和测试后，手写命令会越来越长，而且容易漏依赖。CMake 的价值不是“多一个工具”，而是把目标关系记录下来：库由哪些源文件组成，公开哪些头文件路径，测试链接哪个库。

推荐的构建流程如下：

\begin{lstlisting}[language=bash,caption={独立构建目录}]
cmake -S . -B build -DCMAKE_BUILD_TYPE=Debug
cmake --build build
ctest --test-dir build --output-on-failure
\end{lstlisting}

\code{-S} 指源码目录，\code{-B} 指构建目录。不要把生成文件混进源码目录。Debug 构建适合调试，Release 构建适合性能测量。

\section{头文件和源文件怎么拆}

以成绩统计库为例，头文件只放接口和必要类型：

\begin{lstlisting}[language=C++,caption={summary.hpp 暴露稳定接口}]
#pragma once

#include <optional>
#include <span>

struct ScoreSummary {
    std::size_t count = 0;
    int sum = 0;
    int min_score = 0;
    int max_score = 0;
    double average = 0.0;
};

std::optional<ScoreSummary> summarize_scores(std::span<const int> scores);
\end{lstlisting}

实现文件放算法细节：

\begin{lstlisting}[language=C++,caption={summary.cpp 放实现}]
#include <score/summary.hpp>

#include <algorithm>
#include <numeric>

std::optional<ScoreSummary> summarize_scores(std::span<const int> scores)
{
    if (scores.empty()) {
        return std::nullopt;
    }
    const auto [min_it, max_it] = std::minmax_element(scores.begin(), scores.end());
    const int sum = std::accumulate(scores.begin(), scores.end(), 0);
    return ScoreSummary{scores.size(), sum, *min_it, *max_it,
                        static_cast<double>(sum) / scores.size()};
}
\end{lstlisting}

调用者只需要包含头文件，不应该依赖实现文件里的私有细节。这个边界越干净，后面替换实现越容易。

\section{测试要覆盖失败路径}

很多测试只测普通输入，结果错误路径一直没人碰。成绩统计至少要有这些测试：

\begin{lstlisting}[language=C++,caption={最小测试集合}]
void test_empty_scores()
{
    const std::vector<int> scores;
    assert(!summarize_scores(scores).has_value());
}

void test_single_score()
{
    const std::vector<int> scores{90};
    const auto summary = summarize_scores(scores);
    assert(summary.has_value());
    assert(summary->min_score == 90);
    assert(summary->max_score == 90);
    assert(summary->average == 90.0);
}
\end{lstlisting}

测试函数命名要说清行为，而不是叫 \code{test1}、\code{test2}。失败时，名字就是第一条诊断信息。

\section{调试工具各自解决什么}

编译警告帮你发现可疑代码；单元测试帮你固定行为；调试器帮你观察某次运行；消毒器帮助发现越界、悬空引用和数据竞争。可以给 Debug 构建加上地址消毒器：

\begin{lstlisting}[language=CMake,caption={为测试目标打开地址消毒器}]
target_compile_options(summary_tests PRIVATE -fsanitize=address -fno-omit-frame-pointer)
target_link_options(summary_tests PRIVATE -fsanitize=address)
\end{lstlisting}

这不是替代测试，而是让测试失败得更早、更明确。比如越界访问可能平时“看起来能跑”，在消毒器下会直接报出位置。

\begin{exercisebox}
把成绩统计项目拆成 \filepath{include/score/summary.hpp}、\filepath{src/summary.cpp}、\filepath{tests/summary_tests.cpp}。要求 \code{ctest} 能运行测试；再故意把空输入检查删掉，确认测试能失败。
\end{exercisebox}
